\section{Selection data}
\label{sec:selection}

The main construction uses an exact tensor endofunctor on finite-dimensional
modules whose finite extinction times are unbounded. The example
in~\cite[Section~2]{OAI26findim} is defined over the complex group algebra of a
finitely presented group.

\subsection{Selection data and extinction}
\label{subsec:selection-data}

\begin{definition}
  \label{def:selection-data}
  Let $\kk$ be a field. \emph{Selection data} over $\kk$ consist of a
  $\kk$-algebra $R$ and an $R$-bimodule $\Psi$ on which the two actions of
  $\kk$ agree and which is finitely generated and projective as a right
  $R$-module. The \emph{selection functor} is
  \[
    H=\Psi\otimes_R-\colon\operatorname{mod}_{\mathrm{fd}}R\longrightarrow
    \operatorname{mod}_{\mathrm{fd}}R.
  \]
  The \emph{extinction time} of a
  finite-dimensional left $R$-module $Y$ is the least integer ${t\geq0}$ with
  ${H^tY=0}$, or $\infty$ if there is none. The selection data have
  \emph{unbounded extinction} if the finite extinction times of
  finite-dimensional modules are unbounded.
\end{definition}

The functor $H$ is well defined and exact: $\Psi_R$ is a direct summand of some
$R^n$, so that, as a vector space, $HY$ is a direct summand of $Y^n$, naturally
in $Y$, and exactness is detected on underlying vector spaces; see
\Cref{lemma:selection-matrix} for an explicit description. Iterates of $H$ are
written $H^t$, with $H^0$ the identity.

\begin{example}
  \label{ex:twisted-corner}
  Let ${e\in R}$ be a central idempotent and ${\alpha\colon R\to R}$ a homomorphism
  of unital $\kk$-algebras, and let ${\Psi={}_\alpha(eR)}$ be the bimodule $eR$
  with its right regular action and the left action ${r\cdot x=\alpha(r)x}$.
  Then $\Psi$ is projective as a right module, and ${HY\cong{}_\alpha(eY)}$ is
  the $R$-submodule $eY$ of $Y$ with $r$ acting as $\alpha(r)$: the map
  ${\Psi\otimes_RY\to eY}$, ${x\otimes y\mapsto xy}$, has the inverse
  ${v\mapsto e\otimes v}$, since ${e\otimes xy=ex\otimes y=x\otimes y}$. This is
  the case treated in~\cite[Section~2]{OAI26findim}.
\end{example}

Selection data with unbounded extinction must detect infinitely many
independent numerical invariants of finite-dimensional modules. To make this
precise, let $K_0(R^{\op})$ be the Grothendieck group of finitely generated
projective right $R$-modules. For a finite-dimensional left $R$-module $Y$ the
\emph{rank function} of $Y$ is the additive map
\[
  \chi_Y\colon K_0(R^{\op})\longrightarrow\ZZ,\qquad
  [P]\longmapsto\dim_\kk(P\otimes_RY),
\]
and the bimodule $\Psi$ induces the endomorphism
${\mathsf{T}[P]=[P\otimes_R\Psi]}$
of $K_0(R^{\op})$; here ${P\otimes_R\Psi}$ is finitely generated and projective
since $P$ is a direct summand of some $R^n$. Associativity of the tensor
product gives
\begin{equation}
  \label{eq:rank-shift}
  \chi_Y(\mathsf{T}x)=\chi_{HY}(x)\qquad\text{and}\qquad
  \dim_\kk H^tY=\chi_Y(\mathsf{T}^t[R])
\end{equation}
for every ${x\in K_0(R^{\op})}$ and every ${t\geq0}$. Let $\mathcal{F}$ be the
$\mathbb{Q}$-vector space of functions from the set of isomorphism classes of
finite-dimensional left $R$-modules to $\mathbb{Q}$, and let
\[
  \chi\colon K_0(R^{\op})\otimes_\ZZ\mathbb{Q}\longrightarrow\mathcal{F},
  \qquad x\longmapsto(Y\mapsto\chi_Y(x)).
\]
For the selection data, put
\[
  \mathcal{V}=\chi\bigl(\operatorname{span}_{\mathbb{Q}}
  \set{\mathsf{T}^t[R]}[t\geq0]\bigr)\subseteq\mathcal{F}.
\]

\begin{proposition}
  \label{prop:rank-obstruction}
  Let $(R,\Psi)$ be selection data, and suppose that ${r=\dim_{\mathbb Q}\mathcal
  V<\infty}$. Then every
  finite-dimensional left $R$-module of finite extinction time has extinction
  time at most~$r$.
\end{proposition}

\begin{proof}
  By~\eqref{eq:rank-shift}, if ${\chi(x)=0}$ then ${\chi(\mathsf{T}x)=0}$, since
  $HY$ is again finite-dimensional. Hence, $\mathsf{T}$ induces an endomorphism
  $\overline{\mathsf{T}}$ of $\mathcal{V}$, and $\mathcal{V}$ is spanned by the
  vectors $\overline{\mathsf{T}}^tv$ for ${t\geq0}$, where ${v=\chi([R])}$. The
  Fitting decomposition of $\overline{\mathsf{T}}$ gives
  ${\mathcal{V}=\mathcal{V}_0\oplus\mathcal{V}_1}$, where $\overline{\mathsf{T}}$
  is nilpotent on $\mathcal{V}_0$, hence
  ${\overline{\mathsf{T}}^r\mathcal{V}_0=0}$,
  and invertible on $\mathcal{V}_1$. Write ${v=v_0+v_1}$ accordingly, and let
  ${\mathcal{V}_2\subseteq\mathcal{V}_1}$ be the span of the vectors
  $\overline{\mathsf{T}}^tv_1$ for ${t\geq0}$. The restriction of
  $\overline{\mathsf{T}}$ to $\mathcal{V}_2$ is injective, hence bijective, so
  for every $t_0$ the vectors $\overline{\mathsf{T}}^tv_1$ with ${t\geq t_0}$
  span ${\overline{\mathsf{T}}^{t_0}\mathcal{V}_2=\mathcal{V}_2}$.

  Let $Y$ be a finite-dimensional module of finite extinction time $t_0$, and
  let ${\phi_Y\colon\mathcal{V}\to\mathbb{Q}}$ be evaluation at $Y$. Then
  ${\phi_Y(\overline{\mathsf{T}}^tv)=\dim_\kk H^tY=0}$ for ${t\geq t_0}$
  by~\eqref{eq:rank-shift}. For ${t\geq\max(t_0,r)}$ we have
  ${\overline{\mathsf{T}}^tv=\overline{\mathsf{T}}^tv_1}$, so $\phi_Y$ vanishes on
  $\mathcal{V}_2$. Hence,
  \[
    \dim_\kk H^rY=\phi_Y(\overline{\mathsf{T}}^rv_0)+
    \phi_Y(\overline{\mathsf{T}}^rv_1)=0.\qedhere
  \]
\end{proof}

\begin{corollary}
  \label{coro:rank-obstruction}
  Selection data $(R,\Psi)$ over $\kk$ have bounded extinction in each of the
  following cases.
  \begin{enumerate}
    \item\label{it:rank-finite-list} There are finitely many finitely
      generated projective right $R$-modules such that every finitely
      generated projective right $R$-module is a finite direct sum of copies
      of them; for example, $R$ is finite-dimensional.
    \item\label{it:rank-noetherian} The algebra $R$ is commutative and
      noetherian.
    \item\label{it:rank-localisation} The algebra $R$ is a universal
      localisation of a finite-dimensional hereditary algebra.
  \end{enumerate}
\end{corollary}

\begin{proof}
  By \Cref{prop:rank-obstruction}, it suffices to show that $\chi$ has image
  of finite dimension. In case~\eqref{it:rank-finite-list} the classes of the
  finitely many modules span ${K_0(R^{\op})\otimes_\ZZ\mathbb{Q}}$; if $R$ is
  finite-dimensional, the Krull--Schmidt theorem shows that every finitely
  generated projective module is a finite direct sum of indecomposable direct
  summands of $R$, of which there are finitely many up to isomorphism. In
  case~\eqref{it:rank-localisation}, let ${R=B_\Sigma}$ be a universal
  localisation of a finite-dimensional hereditary algebra $B$. By a theorem of
  Schofield~\cite[Theorem~2.3]{Sch07a}, every universal localisation of $B$ is
  of the form studied in~\cite{Sch07a}, and for these the map from
  $K_0(B^{\op})$ to $K_0(B_\Sigma^{\op})$, for right modules as
  in~\cite{Sch07a}, is surjective~\cite[Lemma~4.1]{Sch07a}. Since
  $K_0(B^{\op})$ is finitely generated by case~\eqref{it:rank-finite-list}, so
  is $K_0(B_\Sigma^{\op})$.

  In case~\eqref{it:rank-noetherian}, for a finitely generated projective
  $R$-module $P$, the function
  \[
    \rho_P\colon\operatorname{Spec}R\longrightarrow\ZZ,\qquad
    \mathfrak{p}\longmapsto
    \dim_{\kappa(\mathfrak{p})}P\otimes_R\kappa(\mathfrak{p})
  \]
  is locally constant~\cite[Tag~00NX, Lemma~10.78.2]{Stacks}. The space
  $\operatorname{Spec}R$ has finitely many irreducible
  components~\cite[Tag~00FR, Lemma~10.31.6]{Stacks}; each of them is connected
  and contained in a connected component, so $\operatorname{Spec}R$ has
  finitely many connected components $Z_1,\dots,Z_c$, and $\rho_P$ takes a
  constant value $\rho_P(Z_i)$ on each of them. A finite-dimensional module $Y$
  has finite length, its composition factors are of the form $R/\mathfrak{m}$
  with $\mathfrak{m}$ maximal and $R/\mathfrak{m}$ finite-dimensional, and
  \[
    \dim_\kk(P\otimes_RR/\mathfrak{m})=\rho_P(\mathfrak{m})\dim_\kk
    R/\mathfrak{m}.
  \]
  Since ${P\otimes_R-}$ is exact,
  \[
    \chi_Y([P])=\sum_{i=1}^c\rho_P(Z_i)\,d_i(Y),
  \]
  where $d_i(Y)$ is the sum of the dimensions of the composition factors
  $R/\mathfrak{m}$ of $Y$ with ${\mathfrak{m}\in Z_i}$. Hence, the image of
  $\chi$ lies in the span of the $c$ functions $d_i$.
\end{proof}

In case~\eqref{it:rank-noetherian}, the argument bounds the dimension of the
image of $\chi$ by the number of connected components, without requiring
$K_0(R)$ to have finite rank. In the example below, $\mathcal V$ is
infinite-dimensional; see \Cref{rem:selection-visible-rank}.

\subsection{The group}
\label{subsec:group}

We write ${[x,y]=xyx^{-1}y^{-1}}$ for elements of a group, and put
${I=\set{1,2,3}}$. Let $\varepsilon_1,\varepsilon_2,\varepsilon_3$ be the
standard basis of $\ZZ^3$, and for ${n=(n_1,n_2,n_3)\in\ZZ^3}$ and commuting
elements $T_1,T_2,T_3$ write ${T^n=T_1^{n_1}T_2^{n_2}T_3^{n_3}}$.

\begin{definition}
  \label{def:selection-group}
  Let $G$ be the group with generators $T_1$, $T_2$, $T_3$ and $U_i(r)$,
  $V_i(r)$, $W_{ij}(r)$ for ${i,j\in I}$ with ${i\neq j}$ and ${r\in\ZZ}$, subject
  to the following relations, for all integers $r$ and $s$ and all indices
  satisfying the stated conditions: the $T_\ell$ commute pairwise;
  \begin{equation}
    \label{eq:group-conjugation}
    T_\ell S(r)T_\ell^{-1}=S(r+\lambda_S(\varepsilon_\ell))
  \end{equation}
  for every generator type ${S\in\set{U_i,V_i,W_{ij}}}$, where the
  \emph{weights} are the homomorphisms ${\lambda_S\colon\ZZ^3\to\ZZ}$ given by
  ${\lambda_{U_i}(n)=-n_i}$, ${\lambda_{V_i}(n)=n_i}$ and
  ${\lambda_{W_{ij}}(n)=n_i-n_j}$;
  \begin{equation}
    \label{eq:group-relations}
    \begin{gathered}
      U_i(r)^2=1,\qquad
      [U_i(r),U_j(s)]=[V_i(r),V_j(s)]=[U_i(r),V_j(s)]=1\quad(i\neq j),\\
      [W_{ij}(s),U_h(r)]=1\quad(h\neq i),\qquad
      [W_{ij}(s),V_h(r)]=1\quad(h\neq j),\\
      [U_i(r),W_{ij}(s)]=U_j(r+s),\qquad[W_{ij}(s),V_j(r)]=V_i(r+s).
    \end{gathered}
  \end{equation}
\end{definition}

No relation between two generators of the same type with different parameters
and no relation between two generators of types $W_{ij}$ is imposed. The main
preprint describes $G$ as motivated by finitely presented matrix groups of
Abels type, and its commutator arguments parallel the proof
of~\cite[Proposition~4.9, relations (4.9) and (4.10)]{San22}; the proofs below
are self-contained.

\begin{proposition}
  \label{prop:group-presentation}
  The group $G$ is finitely presented. More precisely, put ${u_i=U_i(0)}$,
  ${v_i=V_i(0)}$ and ${w_{ij}=W_{ij}(0)}$. Then $G$ is generated by the fifteen
  elements $T_\ell$, $u_i$, $v_i$ and $w_{ij}$, subject to the following
  finitely many relations: the $T_\ell$ commute pairwise; the relations
  \eqref{eq:group-relations} with ${r=s=0}$; and ${[T_h,u_i]=[T_h,v_i]=1}$ for
  ${h\neq i}$, together with ${[T_iT_j,w_{ij}]=[T_k,w_{ij}]=1}$ for
  ${\set{i,j,k}=I}$.
\end{proposition}

\begin{proof}
  Let $P$ be the group with the stated presentation. For a type $S$ write $s_S$
  for the generator $u_i$, $v_i$ or $w_{ij}$, and choose
  ${q_{U_i}=-\varepsilon_i}$, ${q_{V_i}=\varepsilon_i}$ and
  ${q_{W_{ij}}=\varepsilon_i}$, so that ${\lambda_S(q_S)=1}$. The vectors $b$ with
  ${[T^b,s_S]=1}$ among the relations of $P$ form a basis of the kernel $K_S$ of
  $\lambda_S$: for $U_i$ and $V_i$ the kernel is ${\set{n}[n_i=0]}$, with basis
  $\varepsilon_h$ for ${h\neq i}$, and for $W_{ij}$ it is ${\set{n}[n_i=n_j]}$,
  whose elements are ${n_i(\varepsilon_i+\varepsilon_j)+n_k\varepsilon_k}$.
  Hence, $T^b$ commutes with $s_S$ for every ${b\in K_S}$. In $P$ put
  \[
    \widetilde{S}(r)=T^{rq_S}s_ST^{-rq_S}\qquad(r\in\ZZ).
  \]
  For ${n\in\ZZ^3}$ we have ${n-\lambda_S(n)q_S\in K_S}$, and the elements
  $T_1,T_2,T_3$ commute, so
  \begin{equation}
    \label{eq:torus-conjugation}
    T^n\widetilde{S}(r)T^{-n}=\widetilde{S}(r+\lambda_S(n)).
  \end{equation}
  This gives the relations~\eqref{eq:group-conjugation} for the elements
  $\widetilde{S}(r)$, and conjugating ${u_i^2=1}$ by $T^{rq_{U_i}}$ gives
  ${\widetilde{U}_i(r)^2=1}$. Each of the remaining relations
  in~\eqref{eq:group-relations} involves two generators of types $S$ and $S'$
  with prescribed parameters, and it is obtained by conjugating the
  corresponding relation with parameters zero by $T^n$, for an ${n\in\ZZ^3}$
  with the prescribed values of $\lambda_S(n)$ and $\lambda_{S'}(n)$, and using
  \eqref{eq:torus-conjugation}. Such vectors $n$ are, for distinct $i,j,k$:
  \[
    \begin{array}{ll@{\qquad}ll}
      [U_i(r),U_j(s)]\colon&-r\varepsilon_i-s\varepsilon_j, &
      [W_{ij}(s),V_i(r)]\colon&r\varepsilon_i+(r-s)\varepsilon_j,\\{}
      [V_i(r),V_j(s)]\colon&r\varepsilon_i+s\varepsilon_j, &
      [W_{ij}(s),V_k(r)]\colon&s\varepsilon_i+r\varepsilon_k,\\{}
      [U_i(r),V_j(s)]\colon&-r\varepsilon_i+s\varepsilon_j, &
      [U_i(r),W_{ij}(s)]\colon&-r\varepsilon_i-(r+s)\varepsilon_j,\\{}
      [W_{ij}(s),U_j(r)]\colon&(s-r)\varepsilon_i-r\varepsilon_j, &
      [W_{ij}(s),V_j(r)]\colon&(r+s)\varepsilon_i+r\varepsilon_j.\\{}
      [W_{ij}(s),U_k(r)]\colon&s\varepsilon_i-r\varepsilon_k, &&
    \end{array}
  \]
  In the last two
  cases the weight of the third generator is ${\lambda_{U_j}(n)=r+s}$ and
  ${\lambda_{V_i}(n)=r+s}$, as required. Hence, all defining relations of $G$
  hold for the elements $T_\ell$ and $\widetilde{S}(r)$ of $P$.

  Conversely, the relations~\eqref{eq:group-conjugation} give
  ${T^nS(r)T^{-n}=S(r+\lambda_S(n))}$ in $G$ for all ${n\in\ZZ^3}$, which implies
  the relations ${[T^b,S(0)]=1}$ for ${b\in K_S}$; the other relations of $P$ are
  instances of relations of $G$. Thus, there are homomorphisms ${P\to G}$ sending
  $s_S$ to $S(0)$ and ${G\to P}$ sending $S(r)$ to $\widetilde{S}(r)$, both
  fixing the $T_\ell$. Their composites fix all generators, since
  ${\widetilde{S}(0)=s_S}$ and ${T^{rq_S}S(0)T^{-rq_S}=S(r)}$.
\end{proof}

\begin{proposition}
  \label{prop:central-involutions}
  For every ${N\in\ZZ}$ the element
  \[
    z_N=[U_i(a),V_i(b)]\qquad(i\in I,\ a,b\in\ZZ,\ a+b=N)
  \]
  of $G$ does not depend on the choice of $i$, $a$ and $b$. It is central and
  satisfies ${z_N^2=1}$.
\end{proposition}

\begin{proof}
  Let ${i\neq j}$ and ${a,s,b\in\ZZ}$, and put ${x=U_i(a)}$, ${w=W_{ij}(s)}$,
  ${v=V_j(b)}$, ${p}={[x,w]}={U_j(a+s)}$ and ${q}={[w,v]}={V_i(s+b)}$. By the relations
  \eqref{eq:group-relations}, the element $x$ commutes with $v$, and $q$
  commutes with $p$ and with $v$. Using ${xwx^{-1}=pw}$ and ${wvw^{-1}=qv}$, we
  compute
  \[
    xqx^{-1}=x(wvw^{-1}v^{-1})x^{-1}=(pw)v(pw)^{-1}v^{-1}
    =p(qv)p^{-1}v^{-1}=q[p,v].
  \]
  Since $q$ commutes with $p$ and $v$, it commutes with $[p,v]$, and
  multiplying by $q^{-1}$ gives ${[x,q]=[p,v]}$, that is,
  \begin{equation}
    \label{eq:commutator-transfer}
    [U_i(a),V_i(s+b)]=[U_j(a+s),V_j(b)].
  \end{equation}
  If ${a+b=c+d}$ and ${i\neq j}$, then~\eqref{eq:commutator-transfer} with
  ${s=c-a}$ and $b$ replaced by $d$ gives ${[U_i(a),V_i(b)]=[U_j(c),V_j(d)]}$.
  Two splittings at the same index $i$ are both equal to $[U_j(N),V_j(0)]$ for
  some ${j\neq i}$. This proves the independence.

  Let $g$ be a generator of $G$. If $g$ is $U_h(r)$ or $V_h(r)$, write $z_N$ at
  an index ${i\neq h}$; then $g$ commutes with $U_i(a)$ and $V_i(b)$
  by~\eqref{eq:group-relations}. If ${g=W_{ij}(r)}$, write $z_N$ at the index
  ${k\notin\set{i,j}}$, and use the relations ${[W_{ij},U_k]=[W_{ij},V_k]=1}$. If
  ${g=T_\ell}$, then~\eqref{eq:group-conjugation} gives
  \[
    T_\ell z_NT_\ell^{-1}=[U_i(a-\delta_{\ell i}),V_i(b+\delta_{\ell i})]=z_N.
  \]
  Hence, $z_N$ is central. Finally, with ${x=U_i(a)}$ and ${y=V_i(b)}$, centrality
  of ${z_N=[x,y]}$ gives ${xyx^{-1}=z_Ny}$ and then ${x^2yx^{-2}=z_N^2y}$; since
  ${x^2=1}$, we obtain ${z_N^2=1}$.
\end{proof}

\begin{proposition}
  \label{prop:shift-automorphism}
  For every ${c\in\ZZ}$ there is an automorphism $\beta_c$ of $G$ fixing the
  $T_\ell$, the $V_i(r)$ and the $W_{ij}(r)$, with ${\beta_c(U_i(r))=U_i(r+c)}$.
  Moreover, ${\beta_c\beta_d=\beta_{c+d}}$, and ${\alpha_G\coloneqq\beta_1}$
  satisfies ${\alpha_G(z_N)=z_{N+1}}$ for every $N$.
\end{proposition}

\begin{proof}
  The assignment sends each defining relation to a defining relation:
  \begin{itemize}
    \item the relations~\eqref{eq:group-conjugation} for $U_i$ go to
      \[
        T_\ell U_i(r+c)T_\ell^{-1}=U_i(r+c-\delta_{\ell i});
      \]
    \item the relations
      \begin{gather*}
        U_i(r)^2=1,\qquad[U_i(r),U_j(s)]=1,\\
        [U_i(r),V_j(s)]=1,\qquad[W_{ij}(s),U_h(r)]=1
      \end{gather*}
      go to the instances with ${r+c}$ and ${s+c}$ in place of the parameters of
      $U$;
    \item the relation ${[U_i(r),W_{ij}(s)]=U_j(r+s)}$ goes to its instance with
      first parameter ${r+c}$;
    \item all other relations are fixed.
  \end{itemize}
  Hence,
  $\beta_c$ is an endomorphism, and the identities ${\beta_c\beta_d=\beta_{c+d}}$
  and ${\beta_0=\id}$ hold on generators, so $\beta_c$ is an automorphism with
  inverse $\beta_{-c}$. Applying $\beta_1$ to ${z_N=[U_i(a),V_i(b)]}$ gives
  ${[U_i(a+1),V_i(b)]=z_{N+1}}$ by \Cref{prop:central-involutions}.
\end{proof}

\subsection{Finite quotients}
\label{subsec:finite-quotients}

For an integer ${m\geq1}$ let ${S_m=\mathbb{F}_2[t]/(t^m-1)}$. Division by the
monic polynomial ${t^m-1}$ shows that $1,t,\dots,t^{m-1}$ is a basis of $S_m$
over $\mathbb{F}_2$, and $t$ is a unit with inverse $t^{m-1}$; the ring $S_m$
need not be reduced. We write $E_{ab}$ for the matrix units in $M_5(S_m)$ and
$\mathbf{1}$ for the identity matrix.

\begin{proposition}
  \label{prop:finite-quotients}
  For every ${m\geq1}$ there is a homomorphism
  ${\pi_m\colon G\to\operatorname{GL}_5(S_m)}$ with
  \[
    \begin{gathered}
      \pi_m(T_\ell)=D_\ell,\qquad\pi_m(U_i(r))=\mathbf{1}+t^rE_{1,i+1},\\
      \pi_m(V_i(r))=\mathbf{1}+t^rE_{i+1,5},\qquad
      \pi_m(W_{ij}(r))=\mathbf{1}+t^rE_{i+1,j+1},
    \end{gathered}
  \]
  where $D_\ell$ is the diagonal matrix with entry $t$ in position ${\ell+1}$ and
  entries $1$ elsewhere. Its image $F_m$ is finite, and
  ${\pi_m(z_N)=\mathbf{1}+t^NE_{15}}$ for every ${N\in\ZZ}$. The elements
  $\pi_m(z_0),\dots,\pi_m(z_{m-1})$ generate a central subgroup $Z_m$ of $F_m$
  isomorphic to $(\ZZ/2)^m$.
\end{proposition}

\begin{proof}
  Every matrix unit ${A=E_{ab}}$ with ${a\neq b}$ satisfies ${A^2=0}$, so
  ${(\mathbf{1}+xA)^2}={\mathbf{1}+2xA}={\mathbf{1}}$ for ${x\in S_m}$; in particular,
  the matrices above are invertible and the relations ${U_i(r)^2=1}$ hold. The
  diagonal matrices commute, and ${DE_{ab}D^{-1}}={d_ad_b^{-1}E_{ab}}$ for a
  diagonal matrix $D$ with entries $d_a$; for ${D=D_\ell}$ the factors for
  $E_{1,i+1}$, $E_{i+1,5}$ and $E_{i+1,j+1}$ are $t^{-\delta_{\ell i}}$,
  $t^{\delta_{\ell i}}$ and $t^{\delta_{\ell i}-\delta_{\ell j}}$, which gives
  the relations~\eqref{eq:group-conjugation}. If $A$ and $B$ are matrix units
  with ${AB=BA=0}$, then ${\mathbf{1}+xA}$ and ${\mathbf{1}+yB}$ commute; since
  ${E_{ab}E_{cd}=\delta_{bc}E_{ad}}$, this applies to the pairs in the
  commutation relations of~\eqref{eq:group-relations}, by the restrictions on
  the indices: for instance, for $W_{ij}$ and $U_h$ with ${h\neq i}$ the
  products are ${E_{i+1,j+1}E_{1,h+1}=0}$ and
  ${E_{1,h+1}E_{i+1,j+1}}={\delta_{hi}E_{1,j+1}}={0}$. For distinct $a$, $b$ and $c$,
  the matrices ${A=E_{ab}}$ and ${B=E_{bc}}$ satisfy ${BA=0}$ and ${AB=E_{ac}}$, and
  $AB$ has zero product with $A$ and $B$ on both sides; hence
  \[
    [\mathbf{1}+xA,\mathbf{1}+yB]=(\mathbf{1}+xA)(\mathbf{1}+yB)
    (\mathbf{1}+xA)(\mathbf{1}+yB)=\mathbf{1}+xyAB.
  \]
  With ${(a,b,c)=(1,i+1,j+1)}$ and ${(i+1,j+1,5)}$ this gives the last two
  relations of~\eqref{eq:group-relations}, and with ${(1,i+1,5)}$ it gives
  ${\pi_m(z_N)=\mathbf{1}+t^{a+b}E_{15}=\mathbf{1}+t^NE_{15}}$. The image $F_m$ is
  finite since $M_5(S_m)$ is finite.

  The matrix unit $E_{15}$ has zero product on both sides with every matrix unit
  $E_{ab}$, ${a\neq b}$, occurring above, and every $D_\ell$ has equal entries
  in positions $1$ and $5$, so $\pi_m(z_N)$ is central in $F_m$. Since
  ${E_{15}^2=0}$, we have
  \[
    \prod_{q=0}^{m-1}\pi_m(z_q)^{c_q}
    =\mathbf{1}+\Bigl(\sum_{q=0}^{m-1}c_qt^q\Bigr)E_{15}
  \]
  for ${c_q\in\set{0,1}}$, which is the identity only if all $c_q$ vanish.
  Together with ${\pi_m(z_q)^2=\mathbf{1}}$ this gives
  ${Z_m\cong(\ZZ/2)^m}$.
\end{proof}

\begin{corollary}
  \label{coro:z-independent}
  The elements $z_N$, ${N\in\ZZ}$, generate a subgroup of the centre of $G$
  isomorphic to $\bigoplus_{N\in\ZZ}\ZZ/2$. In particular, the centre of $G$ is
  not finitely generated.
\end{corollary}

\begin{proof}
  Let ${N_1<\dots<N_a}$ be integers and choose ${m>N_a-N_1}$. The residues of the
  $N_b$ modulo $m$ are distinct, so \Cref{prop:finite-quotients} shows that
  every nonempty product of the $z_{N_b}$ has nontrivial image in $F_m$. Since
  the $z_N$ are central involutions by \Cref{prop:central-involutions}, they
  generate the stated subgroup. A subgroup of a finitely generated abelian
  group is finitely generated, which this subgroup is not.
\end{proof}

The group $G$ is finitely presented, has an automorphism shifting the
independent central involutions, and has finite quotients $F_m$ separating every
finite subfamily of them.

\subsection{Selection data with unbounded extinction}
\label{subsec:unbounded-extinction}

Let ${R=\mathbb{C}G}$ be the complex group algebra of the group $G$ of
\Cref{def:selection-group}, let ${\alpha\colon R\to R}$ be the linear extension
of $\alpha_G$, and put
\[
  e=\tfrac12(1+z_0)\in R,\qquad\Psi={}_\alpha(eR).
\]

\begin{lemma}
  \label{lemma:selection-iterates}
  Let $R$ be a $\kk$-algebra, ${e\in R}$ a central idempotent and $\alpha$ an
  automorphism of $R$, and let $H$ be the selection functor of
  ${\Psi={}_\alpha(eR)}$. Put ${e_q=\alpha^q(e)}$ and ${p_j=e_0e_1\cdots e_{j-1}}$,
  with ${p_0=1}$. Then for every finite-dimensional left $R$-module $Y$ and every
  ${j\geq0}$ there is an isomorphism
  \[
    H^jY\cong{}_{\alpha^j}(p_jY),
  \]
  natural in $Y$, where ${r\in R}$ acts on the subspace $p_jY$ of $Y$ as
  $\alpha^j(r)$.
\end{lemma}

\begin{proof}
  The elements $e_q$ are central idempotents, since $\alpha$ is an
  automorphism. For ${j=0}$ there is nothing to prove. If
  ${H^jY\cong{}_{\alpha^j}(p_jY)}$, then by \Cref{ex:twisted-corner} the module
  $H^{j+1}Y$ is the image of $e$ acting on ${}_{\alpha^j}(p_jY)$, that is, the
  subspace ${\alpha^j(e)p_jY=p_{j+1}Y}$, with $r$ acting as
  ${\alpha^j(\alpha(r))=\alpha^{j+1}(r)}$. The identifications are restrictions
  of the identity of $Y$, hence natural.
\end{proof}

\begin{theorem}
  \label{thm:selection-unbounded}
  The pair $(R,\Psi)$ is selection data over $\mathbb{C}$ with $R$ finitely
  presented. For every integer ${m\geq1}$ there is a nonzero finite-dimensional
  left $R$-module $Y_m$ of extinction time exactly $m$. In particular,
  $(R,\Psi)$ has unbounded extinction.
\end{theorem}

\begin{proof}
  By \Cref{prop:group-presentation}, the group $G$ has a presentation with
  fifteen generators $g_1,\dots,g_{15}$ and finitely many relators
  $w_1,\dots,w_b$. The algebra with generators $x_1,y_1,\dots,x_{15},y_{15}$
  and relations ${x_iy_i=y_ix_i=1}$ and ${w_c(x,y)=1}$, in which $g_i^{-1}$ is
  replaced by $y_i$, is isomorphic to $R$: the homomorphism to $R$ has
  an inverse given by the linear extension of the group homomorphism from $G$
  to its group of units sending $g_i$ to $x_i$. Hence, $R$ is finitely
  presented. The map $\alpha$ is an automorphism with inverse the linear
  extension of $\alpha_G^{-1}$, and $e$ is a central idempotent by
  \Cref{prop:central-involutions}; so $(R,\Psi)$ is selection data, by
  \Cref{ex:twisted-corner}. Moreover,
  ${e_q=\alpha^q(e)=\frac12(1+z_q)}$ by \Cref{prop:shift-automorphism}.

  Fix ${m\geq1}$, and let $F_m$ and $Z_m$ be as in
  \Cref{prop:finite-quotients}, with ${\overline{z}_q=\pi_m(z_q)}$. Since
  ${Z_m\cong(\ZZ/2)^m}$ with basis $\overline{z}_0,\dots,\overline{z}_{m-1}$,
  there is a character ${\chi_m\colon Z_m\to\mathbb{C}^\times}$ with
  ${\chi_m(\overline{z}_q)=1}$ for ${0\leq q<m-1}$ and
  ${\chi_m(\overline{z}_{m-1})=-1}$.
  Put
  \[
    p_{\chi_m}=\frac{1}{2^m}\sum_{z\in Z_m}\chi_m(z)^{-1}z\in\mathbb{C}F_m,
    \qquad Y_m=(\mathbb{C}F_m)\,p_{\chi_m},
  \]
  regarded as a left $R$-module through $\pi_m$. The coefficient of the
  identity in $p_{\chi_m}$ is $2^{-m}$, so ${Y_m\neq0}$, and $Y_m$ is
  finite-dimensional. For ${h\in Z_m}$, the substitution ${u=hz}$ gives
  ${hp_{\chi_m}=\chi_m(h)p_{\chi_m}}$, and since $Z_m$ is central in $F_m$,
  every ${h\in Z_m}$ acts on $Y_m$ as the scalar $\chi_m(h)$. Hence, $e_q$ acts on
  $Y_m$ as the identity for ${0\leq q<m-1}$ and as zero for ${q=m-1}$. By
  \Cref{lemma:selection-iterates}, $H^jY_m$ is the subspace $p_jY_m$, which is
  ${Y_m\neq0}$ for ${j<m}$ and zero for ${j\geq m}$, since then $p_j$ contains the
  factor $e_{m-1}$.
\end{proof}

\begin{remark}
  \label{rem:selection-visible-rank}
  In the notation of \Cref{prop:rank-obstruction}, we have
  ${\chi_Y(\mathsf{T}^j[R])=\dim p_jY}$ by \Cref{lemma:selection-iterates},
  and the values of these functions on the
  modules $Y_1,\dots,Y_s$ form a triangular matrix with nonzero diagonal, so
  $\mathcal V$ is infinite-dimensional. The module $Y_m$ is the module
  induced from the character $\chi_m$ of $Z_m$, of dimension $[F_m:Z_m]$. The
  integer $m$ enters only through the
  choice of the finite quotient $F_m$ and of the values of $\chi_m$ on the
  central involutions; $R$, $e$ and $\alpha$ are independent of $m$, as required
  in \Cref{sec:simulation}.
\end{remark}
